\documentclass[12pt]{extarticle}
\usepackage{array}
\usepackage[T1]{fontenc}
\usepackage[Russian]{babel}
\usepackage{booktabs}
\usepackage{graphicx}
\usepackage{spreadtab}
\usepackage{numprint}
\usepackage[margin=3cm]{geometry}
\usepackage{float}
\usepackage{xfp}
\usepackage{amsmath} 
\usepackage{indentfirst}  
\begin{document}
\author{Осипов Александр Б06-604}
\title{Лабораторная работа 1.1.4. ИЗУЧЕНИЕ СТАТИСТИЧЕСКИХ ЗАКОНОМЕРНОСТЕЙ НА 
ПРИМЕРЕ ИЗМЕРЕНИЯ ФОНА КОСМИЧЕСКОГО ИЗЛУЧЕНИЯ}
\maketitle

Работа - https://mipt.ru/upload/Кафедра%20общей%20физики/1/lab/lab_114-2023.pdf

\section{Цель}

На примере статистики регистрации фоновых космических частиц изучить статистические закономерности однородного во времени случайного процесса;
проверить возможность описания исследуемого процесса статистическими законами
Пуассона и Гаусса; измерить среднее число регистрируемых космических лучей в секунду и определить погрешность результата.
\vspace{5mm}

\section{Оборудование}

Счётчик Гейгера—Мюллера, компьютер с интерфейсом для
связи со счётчиком, расчётная программа.
\vspace{5mm}

\section{Теоретическая часть (приницип работаы счётчика)}

Счётчик представляет собой наполненный газом
металлический цилиндр с двумя электродами. Одним из электродов (като-
дом) служит сам корпус. Другим (анодом) является тонкая нить, натянутая
вдоль оси цилиндрического корпуса. Необходимое напряжение (400 В) пода-
ётся на счётчик от смонтированного вместе с ним блока питания через повы-
шающий трансформатор. Радиоактивное излучение (космические частицы)
ионизует молекулы газа, которым наполнен счётчик, а также выбивает элек-
троны из его стенок. Образовавшиеся электроны, двигаясь в сильном элек-
трическом ноле между электродами счётчика, соударяются с молекулами
газа, выбивая из них новые — вторичные электроны. Ускоряясь полем, пер-
вичный и вторичные электроны снова ионизуют газ, и т.д. В результате об-
разуется лавина электронов, и через счётчик протекает кратковременный им-
пульс тока (разряд). Этот импульс и регистрируется электрической цепью
установки, оцифровывается платой аналогово-цифрового преобразователя, и
информация о нём через USB-интерфейс передаётся на компьютер.
\vspace{5mm}

\section{Ход работы}

Измерения проводятся с помощью программы на компьютере.

\begin{enumerate}
    \item В разделе "Эксперимент" запускается эксперимент длительностью 4000с.
    \item В процессе сбора данных программой, выполняется демонстрационное задание:
    \begin{enumerate}
        \item Запускается симуляцию с настройками по умолчанию («Распределение 
Пуассона», интенсивность — 1,0, ускорение — 10). Результаты сохраняются в виде текстовых
и графических файлов.
        \item Симуляция запускается в максимально ускоренном режиме (ускорение = 4000),
изменяя среднюю интенсивность числа частиц.
    \end{enumerate}
    \item Данные обрабатываются в программе, чтобы определить статичестические параметры распределения.
\end{enumerate} 

\section{Обработка результатов}

1. Расчет частот разиличных значений для разных временных интервалов (п.14):

\vspace{5mm}
2. Гистограммы распределений (п.15):

\vspace{5mm}
3. Расчет параметров распределения для различных временных интервалов (п.16):

\vspace{5mm}
4. Сравнение полученного распределения с распределением Гаусса и Парето (п.17):

\vspace{5mm}
5. (п.18)

\vspace{5mm}
6. (п.19) 

\end{document}
